% generated by roam.py handbook export -- не редактировать, править заметки
% 20260917210321-характеристика_выпуклого_множества_через_выпуклые_комбинации.org (body aa80ca90b346)

\begin{theorem}[{Характеристика выпуклого множества через выпуклые комбинации}]\label[theorem]{note:42DFF6DC-CC9D-40AA-9D6F-218DCF613833}
  Пусть $V$~--- векторное пространство. $C \subset V$.
  $C$~--- \hyperref[note:316B5BFA-8B9F-4267-9086-D6593AAB5F9B]{выпуклое} тогда и только тогда, когда оно содержит все свои \hyperref[note:8583F986-A6E8-4BBB-B2FB-BA0E9BBE58BF]{выпуклые комбинации}.
\end{theorem}

\begin{proof}
  \proofpart{$\Leftarrow$ (Содержит комбинации, следовательно выпуклое)}
  Если $\lambda \in [0;1]$, то частным случаем ($n=2, \lambda_1=\lambda, \lambda_2 = 1-\lambda, x_1=x, x_2 = y$) \hyperref[note:8583F986-A6E8-4BBB-B2FB-BA0E9BBE58BF]{выпуклой комбинации} является $\lambda x + (1-\lambda)y\ \forall x, y \in C$.
  А это в точности определение \hyperref[note:316B5BFA-8B9F-4267-9086-D6593AAB5F9B]{выпуклого множества}.

  \proofpart{$\Rightarrow$ (Выпуклое, следовательно содержит комбинации)}
  Строим индукцию.
  База это определение \hyperref[note:316B5BFA-8B9F-4267-9086-D6593AAB5F9B]{выпуклого множества}.

  Положим, что для $\forall n \le m$ утверждение  из заголовка верно.
  Увеличим размерность суммы.

  Давайте запишем произвольную \hyperref[note:8583F986-A6E8-4BBB-B2FB-BA0E9BBE58BF]{выпуклую комбинацию}
  $\sum\limits_{k=1}^{m+1} \lambda_k x_k,\ x_k \in C, \sum\limits_{k=1}^{m+1} \lambda_k = 1$. Если $\lambda_{m+1} = 1$, то остальные $\lambda_k = 0$ и комбинация равна $x_{m+1} \in C$. Далее считаем $\lambda_{m+1} \ne 1$.

  Тогда:
  $\sum\limits_{k=1}^{m+1} \lambda_k x_k = \sum\limits_{k=1}^{m} \lambda_k x_k + \lambda_{m+1}x_{m+1}$

  Пусть $y = \sum\limits_{k=1}^m \frac{\lambda_k x_k}{1 - \lambda_{m+1}}$. Но $\sum\limits_{k=1}^m \frac{\lambda_k}{1-\lambda_{m+1}} = \frac{1-\lambda_{m+1}}{1-\lambda_{m+1}} = 1$, и коэффициенты неотрицательны ($1 - \lambda_{m+1} > 0$), следовательно $y$~--- \hyperref[note:8583F986-A6E8-4BBB-B2FB-BA0E9BBE58BF]{выпуклая комбинация} и можем применить индукционное предположение для размерности $m$.
  Тогда $y \in C$.

  Получили, что $\sum\limits_{k=1}^{m+1} \lambda_k x_k = \sum\limits_{k=1}^{m} \lambda_k x_k + \lambda_{m+1}x_{m+1} = (1-\lambda_{m+1})y + \lambda_{m+1}x_{m+1}$
  А по \hyperref[note:316B5BFA-8B9F-4267-9086-D6593AAB5F9B]{определению выпуклости} этот объект лежит в множестве.
\end{proof}
